\section{Asking for Less}
\label{sec:ch03-asking-for-less}

\index{projection!columns nobody asked for}%
Read statement one again. The client asked for \code{title} and the speaker's
\code{name}; SQLite was asked for six columns, including an abstract that can
run to a paragraph. Statement count was the interesting number because it grows
with the list. Column count does not grow, which is why it is second, but it is
paid on every row of every request forever.

\index{QueryContext@\code{QueryContext<T>}}%
\index{selection set!definition}%
The fix is a parameter. \code{QueryContext<Session>} is Hot Chocolate handing
the resolver the client's selection set, which is the set of fields inside the
braces the client wrote, in a form \code{IQueryable} understands, and
\code{With} applies it:

\begin{minted}[highlightlines={1,14,18,27,29}]{csharp}
using GreenDonut.Data;
using Microsoft.EntityFrameworkCore;

namespace Sessions;

[QueryType]
public static class Query
{
    /// <summary>
    /// Every session on the schedule, earliest first.
    /// </summary>
    public static async Task<List<Session>> GetSessionsAsync(
        ConferenceContext db,
        QueryContext<Session> query,
        CancellationToken ct)
        => await db.Sessions
            .OrderBy(s => s.StartsAt)
            .With(query)
            .ToListAsync(ct);

    /// <summary>
    /// One session by its identifier, or null if no session has it.
    /// </summary>
    public static async Task<Session?> GetSessionByIdAsync(
        int id,
        ConferenceContext db,
        QueryContext<Session> query,
        CancellationToken ct)
        => await db.Sessions
            .Where(s => s.Id == id)
            .With(query)
            .FirstOrDefaultAsync(ct);
}
\end{minted}

Both fields take the parameter, and \code{GetSessionByIdAsync} had to be
rewritten to get it: a \code{QueryContext} applies to a query, so the
predicate moves out of \code{FirstOrDefaultAsync} into a \code{Where} and the
projection goes between them. That matters more than it looks, and
this section comes back to it.

\code{AddQueryContext} in \code{Program.cs} is what makes that parameter
resolvable, which is why it was there in
section~\ref{sec:ch03-the-conference-in-sqlite}. The order of the two calls
before \code{ToListAsync} matters and the failure is unpleasant: \code{With}
appends a \code{Select} that rebuilds each \code{Session} from the fields that
were asked for, so sorting after it sorts by a property of a half-built object
and EF Core gives up, quoting back a constructor call you never wrote. Sort
first, then project.

\index{Session@\code{Session}!is a positional record}%
There is an older attribute for this, \code{[UseProjection]}, and it cannot be
used in this book at all. Its provider builds the projected row with
\code{Expression.New} followed by member assignment, so it needs a type with a
parameterless constructor and settable properties. \code{Session} is a
positional record, and every request against a field wearing that attribute
fails at run time with \code{Type 'Session' does not have a default
constructor}. Not at build time, and not at schema build: on the first request.
The attribute has behaved this way for as long as this book's two versions have
existed, so it is a property of how it projects rather than something to wait
out. If you are reading a tutorial that uses it, this is why it is not here.

Send the request again and statement one has shrunk:

\begin{minted}{text}
-- statement 1
SELECT "s"."Title"
FROM "Sessions" AS "s"
ORDER BY "s"."StartsAt"
\end{minted}

\index{projection!starves a resolver|(}%
And every session's speaker is now \code{null}. The server sends this on one
line; it is broken across four here to fit the measure, and nothing else about
it is changed.

\begin{minted}[fontsize=\footnotesize]{json}
{"data":{"sessions":[
  {"title":"Schemas That Outlive Their Authors","speaker":null},
  {"title":"Reading a Query Plan","speaker":null},
  {"title":"Paging Without Offsets","speaker":null},
  {"title":"The Bank That Split Its Graph","speaker":null}]}}
\end{minted}

Status 200. No error key. Nothing in the log, nothing on the console, and a
response that is the right shape and wrong. This is the worst failure in the
chapter and an optimization that worked introduced it.

\index{Parent@\code{[Parent]}!declaring what a resolver reads}%
The cause is in statement two, which asked for one key rather than three.
\code{SpeakerId} is not in the client's selection set, because
\code{[GraphQLIgnore]} keeps it out of the schema entirely, so the projection
did not select it, so every \code{Session} the resolver received had
\code{SpeakerId} left at its default of zero, and no speaker has id zero. The
projection knows what the client asked for. It does not know what your
resolvers read.

So tell it. \code{[Parent]} takes an argument for exactly this:

\begin{minted}[highlightlines={10}]{csharp}
namespace Sessions;

[ObjectType<Session>]
public static partial class SessionType
{
    /// <summary>
    /// The person giving this session.
    /// </summary>
    public static async Task<Speaker?> GetSpeakerAsync(
        [Parent(nameof(Session.SpeakerId))] Session session,
        ISpeakerByIdDataLoader speakerById,
        CancellationToken ct)
        => await speakerById.LoadAsync(session.SpeakerId, ct);
}
\end{minted}

\begin{minted}{text}
-- statement 1
SELECT "s"."Title", "s"."SpeakerId"
FROM "Sessions" AS "s"
ORDER BY "s"."StartsAt"
\end{minted}

Four columns saved, correct answers, still two statements.

\index{projection!starves a resolver|)}%
The rule to take away is a dependency rule, not a syntax one. A projection
narrows the parent to what the client named, and a resolver that reads anything
else has to declare it or it will read a default. I would rather this failed
loudly, and it does not, so treat \code{[Parent(...)]} as part of writing any
resolver that touches the parent's data.

\index{projection!blast radius}%
Note also which field broke. \code{sessionById} was starved in exactly the same
way and for the same reason, because it took a \code{QueryContext} of its own
in the listing above. The defect follows the projection rather than the field,
so adding one to a query is a change to every resolver on the type it returns,
including ones somebody else wrote.

\index{paging!cursor}%
Paging is the last of the three, and after the first two it is almost dull.
\code{sessions} returns every row in the table and always has. One attribute
changes that, and it also changes the shape of the field, so it is the largest
edit in the chapter and the smallest amount of new code:

\begin{minted}[highlightlines={10,12-13,15-17}]{csharp}
using GreenDonut.Data;
using Microsoft.EntityFrameworkCore;

namespace Sessions;

[QueryType]
public static class Query
{
    /// <summary>
    /// Sessions on the schedule, earliest first.
    /// </summary>
    [UsePaging(DefaultPageSize = 2, MaxPageSize = 10)]
    public static IQueryable<Session> GetSessions(
        ConferenceContext db,
        QueryContext<Session> query)
        => db.Sessions
            .OrderBy(s => s.StartsAt)
            .With(query);

    /// <summary>
    /// One session by its identifier, or null if no session has it.
    /// </summary>
    public static async Task<Session?> GetSessionByIdAsync(
        int id,
        ConferenceContext db,
        QueryContext<Session> query,
        CancellationToken ct)
        => await db.Sessions
            .Where(s => s.Id == id)
            .With(query)
            .FirstOrDefaultAsync(ct);
}
\end{minted}

\index{IQueryable@\code{IQueryable}!returned unexecuted}%
The method now returns the query rather than its results, unexecuted and not
awaited, because the paging middleware has to add the slice before anybody runs
it. Losing the \code{async} and the \code{CancellationToken} is part of that:
the middleware awaits on your behalf. \code{GetSessionByIdAsync} keeps both,
because nothing was added to it.

The summary line changed too, from \enquote{Every session on the schedule} to
\enquote{Sessions on the schedule}, and that is not tidying. It is now a lie if
it says every: the field returns a page. A doc comment on a resolver is a
schema description, so this one is API copy, and
chapter~\ref{ch:hot-chocolate-zero} is where that stopped being surprising.

A default page of two is absurd for a real service and useful for a book with
four sessions in it, since it makes the second page exist.

The schema gains a connection, with its types named after the field. The
exported file carries a description on every field and argument below, plus the
\code{@listSize} and \code{@cost} directives on \code{sessions}; the
definitions are quoted here without them, because the shape is the point and
chapter~\ref{ch:hot-chocolate-zero} already showed what the descriptions look
like.

\begin{graphqlsdl}
sessions(
  first: Int
  after: String
  last: Int
  before: String
): SessionsConnection

type SessionsConnection {
  pageInfo: PageInfo!
  edges: [SessionsEdge!]
  nodes: [Session!]
}

type SessionsEdge {
  cursor: String!
  node: Session!
}
\end{graphqlsdl}

\index{connection!definition}%
\index{Relay!cursor connections specification}%
That shape is a connection, and it comes from Relay, Facebook's GraphQL
client, whose cursor connections specification asks for two fields on the
connection and two on each edge~\autocite{relayconnections}. A connection
wraps the page rather than returning it: \code{edges} holds the items, each
paired with the cursor that names its position so that a client can ask for
the page after any of them, and \code{pageInfo} says whether there is more.
\code{nodes} is not in that specification. It is Hot Chocolate's shortcut for
a client that wants the items and has no use for their cursors, and it is the
one every request in this book reads.

\index{breaking change!introduced by paging}%
Every query written against this service before now is broken. \code{sessions
\{ title \}} is no longer legal, because \code{sessions} is not a list any
more, and the replacement is \code{sessions \{ nodes \{ title \} \}}. That is
fine here, where the only client is you and the only history is two chapters.
It is the single most expensive kind of change to make once the graph is
federated and other teams' clients are on it, and
chapter~\ref{ch:breaking-changes} is about catching exactly this before it
ships rather than after.

\begin{minted}{text}
POST /graphql HTTP/1.1
Host: localhost:5001
Content-Type: application/json

{"query":"{ sessions { nodes { title } pageInfo { hasNextPage } } }"}
\end{minted}

Two sessions come back and \code{hasNextPage} is true, because no \code{first}
was given and \code{DefaultPageSize} decided.

\begin{minted}{text}
-- statement 1
SELECT "s"."Title"
FROM "Sessions" AS "s"
ORDER BY "s"."StartsAt"
LIMIT @p
\end{minted}

One column this time, because nothing asked for a speaker and so nothing
needed \code{SpeakerId}. The projection and the page compose without knowing
about each other.

\index{paging!LIMIT pushed into SQL}%
\code{LIMIT} is in the statement, which is the thing to check and the thing
that is easy to get wrong. Paging that happens after the rows arrive has read
the whole table to show you two rows of it, and it looks identical from
outside. The way to know is to read the statement.

\begin{onhc14}
All of this chapter compiles unchanged on Hot Chocolate~14 except the
projection. \code{[DataLoader]} is there, the generated
\code{ISpeakerByIdDataLoader} is there and registered the same unasked way,
\code{[UsePaging]} produces the same \code{SessionsConnection}, and even
\code{[Parent(nameof(Session.SpeakerId))]} compiles.

\code{QueryContext<Session>} also resolves, which makes the failure confusing.
What is missing is the \code{IQueryable} overload of \code{With}, so the
compiler reports the only overload it can see, which takes a DataLoader:

\begin{minted}[fontsize=\footnotesize, breakanywhere]{text}
error CS0411: The type arguments for method
'GreenDonutQueryContextDataLoaderExtensions.With<TKey, TValue>(
IDataLoader<TKey, TValue>, QueryContext<TValue>?)' cannot be inferred from
the usage. Try specifying the type arguments explicitly.
\end{minted}

So on 14 the sessions query selects every column and the rest of the chapter
stands. The counts are the ones this chapter ends on: two statements for a page
of sessions with their speakers, one without. Batching is not among the things
you are missing.

Two further differences you will see in the log belong to Entity Framework
Core rather than to Hot Chocolate, and only appear because 14 does not target
.NET~10 and so runs an older EF Core. Parameters are named \code{@\_\_p\_0}
and \code{@\_\_ids\_0} rather than \code{@p} and \code{@ids1}, and the
DataLoader's key list arrives as a \code{json\_each} subquery rather than a
plain \code{IN} list. Appendix~\ref{app:hc14} carries the full mapping.
\end{onhc14}
